% Independent mathematical derivation. Prefix CPP. Uses the proved CTP receiver.
\section{Every Cayley coefficient as a convergent original-prime-power trace}
\label{sec:cayley-prime-power-formula}

\paragraph{Source and unchanged objects.}
Retain exactly $H(Z)=\xi(1/2+iZ/2)/8$, the parameter $b$,
the full root multiplicities, and $c_k(b)$ of (CTP1)--(CTP18).
The arithmetic decomposition is that used by Enrico Bombieri
and Jeffrey C. Lagarias,
\href{https://doi.org/10.1006/jnth.1999.2392}{\emph{Complements to
Li's criterion for the Riemann hypothesis}}, and by Jeffrey C.
Lagarias, \href{https://arxiv.org/abs/math/0404394v4}{\emph{Li
Coefficients for Automorphic $L$-Functions}}, Section 4, original
author source \path{math-0404394v4.tex}, lines 1425--1604,
especially equations \texttt{N405b}, \texttt{N405e} and Lemma
\texttt{le42}. That source range was read directly. The formulas
below derive the arithmetic expression for the present parameter
$b>1$ and its exact Cayley receiver; they are not an identification
of the full support spectrum with the ordinary prime list.

Fix $b>1$ and put $q=(1+b)/2>1$. Write
$L_\xi(s)=\xi'(s)/\xi(s)$ and use $\psi=\Gamma'/\Gamma$.
The functional equation $\xi(1-s)=\xi(s)$ gives
$L_\xi^{(r)}(1-q)=(-1)^{r+1}L_\xi^{(r)}(q)$. The original
coordinate therefore gives
\begin{equation}\tag{CPP1}
 \ell^{(r)}(ib)=(-i/2)^{r+1}L_\xi^{(r)}(q),\qquad
 c_0(b)=\frac{L_\xi(q)}{2b},\qquad
 c_k(b)=\frac14\sum_{r=1}^k\binom{k-1}{r-1}
                   \frac{b^{r-1}}{r!}L_\xi^{(r)}(q)
 \quad(k\ge1).
\end{equation}
Here the first identity includes $r=0$. To check the last sign
and constant, substitute it in (CTP18); the scalar product
$-(2i)^{r-1}(-i)^{r+1}/2^{r+1}$ is exactly $1/4$.

The defining factors of $\xi$ and the ordinary Euler product of
the actual $\zeta$ give the absolutely convergent expression
\begin{equation}\tag{CPP2}
 L_\xi(q)=\frac1q+\frac1{q-1}-\frac{\log\pi}{2}
             +\frac12\psi(q/2)
             -\sum_p\sum_{l\ge1}(\log p)p^{-lq}.
\end{equation}
For completeness, unique prime factorization expands
$\prod_p(1-p^{-s})^{-1}$ into the absolutely convergent Dirichlet
series for $\zeta(s)$ when $\Re s>1$. Its logarithm is
$\sum_{p,l}p^{-ls}/l$, normally convergent on smaller closed
half-planes. Termwise differentiation yields the last term in
(CPP2), and each further derivative is justified by domination
by a convergent series $\sum_{n\ge2}(\log n)^M n^{-q'}$ with
$1<q'<q$. Thus for $r\ge1$,
\begin{equation}\tag{CPP3}
\begin{split}
 L_\xi^{(r)}(q)
  ={}&(-1)^r r!\{q^{-r-1}+(q-1)^{-r-1}\}
       +2^{-r-1}\psi^{(r)}(q/2)\\
    &+(-1)^{r+1}\sum_p\sum_{l\ge1}
                    (\log p)(l\log p)^r p^{-lq}.
\end{split}
\end{equation}
Every prime power is present, with its original logarithmic
weight, exponent, and sign. Neither $q=1$ nor a limiting
interpretation of this prime series is used.

Define explicitly the finite polynomial
\begin{equation}\tag{CPP4}
 L_{k-1}^{(1)}(X)=\sum_{j=0}^{k-1}\binom{k}{j+1}
                         \frac{(-X)^j}{j!}\quad(k\ge1).
\end{equation}
It is the associated Laguerre polynomial in the convention
specified by this formula. The identity
$\binom{k-1}{r-1}/r!=\binom{k}{r}/(k(r-1)!)$
then gives
\[
 \sum_{r=1}^k\binom{k-1}{r-1}
       \frac{(-1)^{r+1}b^{r-1}x^r}{r!}
           =\frac{x}{k}L_{k-1}^{(1)}(bx).
\]
Consequently the completely evaluated arithmetic formula for
each $k\ge1$ is
\begin{equation}\tag{CPP5}
\begin{split}
 4c_k(b)={}&
  -\frac1{q^2}\left(1-\frac bq\right)^{k-1}
  -\frac1{(q-1)^2}\left(1-\frac b{q-1}\right)^{k-1}\\
 &+\sum_{n=0}^\infty\frac1{(q+2n)^2}
                         \left(1-\frac b{q+2n}\right)^{k-1}\\
 &+\frac1k\sum_p\sum_{l\ge1}
       l(\log p)^2p^{-lq}L_{k-1}^{(1)}(bl\log p).
\end{split}
\end{equation}
The two displayed rational terms come from the original factors
$s$ and $s-1$ of $\xi$. The gamma term is obtained from
$\psi^{(r)}(t)=(-1)^{r+1}r!\sum_{n\ge0}(t+n)^{-r-1}$ for
$r\ge1,t>0$. One proof is to differentiate the logarithm of
Euler's limit
$\Gamma(t)=\lim_{N\to\infty}N!N^t/[t(t+1)\cdots(t+N)]$:
after the first derivative, all subsequent derivatives have
normally convergent series on compact subsets of $t>0$.
Insertion in the finite sum in (CPP1) gives the displayed
gamma series by the binomial theorem. For fixed $k$ that series
converges absolutely, as its summands are $O(n^{-2})$.
The prime series converges absolutely since its summands in
absolute value are bounded by a constant times
$(\log(p^l))^{k+1}(p^l)^{-q}$, and prime powers form a subset
of the integers with no repetitions among distinct pairs $(p,l)$.

The exact endpoint cancellation can also be recorded without
discarding its source. The $n=0$ gamma term in (CPP5) is the
negative of its first rational term. Define the three retained
quantities
\begin{equation}\tag{CPP6}
\begin{split}
 A_k(b)&=\sum_{n=1}^\infty(q+2n)^{-2}
                          (1-b/(q+2n))^{k-1},\\
 E_k(b)&=(q-1)^{-2}(1-b/(q-1))^{k-1},\\
 P_k(b)&=\frac1k\sum_p\sum_{l\ge1}
       l(\log p)^2p^{-lq}L_{k-1}^{(1)}(bl\log p).
\end{split}
\end{equation}
The exact identity is $4c_k=A_k-E_k+P_k$; the two canceled
terms remain explicitly available in (CPP5). Moreover
$0<b/(q+2n)<2$ for $n\ge1$, so
\begin{equation}\tag{CPP7}
 |A_k(b)|\le\sum_{n=1}^\infty(q+2n)^{-2}
     \le(q+2)^{-2}+\frac1{2(q+2)}\quad(k\ge1).
\end{equation}
The second inequality bounds the decreasing tail by its
integral. Hence the complete zero detector (CTP24) is also
equivalent to boundedness of the one precise prime discrepancy
$\{P_k(b)-E_k(b):k\ge1\}$: all real roots give
$|c_k|\le c_0$ and then (CPP7) gives boundedness; conversely
bounded discrepancy makes $c_k$ bounded, whereas any off-real
root makes its limsup exponential rate $R_b>1$ by (CTP22).
The exact sharper inequalities are
$|A_k(b)-E_k(b)+P_k(b)|\le4c_0(b)$ for every $k$.
Neither their boundedness nor their validity for all $k$ is
asserted here. The formula exhibits the required cancellation
against $E_k$, whose factor is
$1-b/(q-1)=-(b+1)/(b-1)$; a bound on one retained term alone
does not evaluate the discrepancy.

\subsection{Explicit finite tails for each original prime sum}

For an integer $m\ge0$ and $N>1$, set
\begin{equation}\tag{CPP8}
 I_m(N,q)=N^{1-q}\sum_{j=0}^m
       \frac{m!}{(m-j)!}\frac{(\log N)^{m-j}}{(q-1)^{j+1}}
       =\int_N^\infty(\log x)^m x^{-q}dx.
\end{equation}
Repeated integration by parts proves the equality, the endpoint
at infinity vanishing because $q>1$. If
$N\ge\exp(m/q)$, the integrand is decreasing on $[N,\infty)$,
as its derivative has sign $m-q\log x$.
Let $\Lambda(n)=\log p$ when $n=p^l$ and zero otherwise;
this symbol is the von Mangoldt function, distinct from the
Li coefficients $\Lambda_n(b)$. Since $0\le\Lambda(n)\le\log n$,
for integer $N\ge\exp((r+1)/q)$ one has
\begin{equation}\tag{CPP9}
 \sum_{p^l>N}(\log p)(l\log p)^r p^{-lq}
 =\sum_{n>N}\Lambda(n)(\log n)^r n^{-q}
 \le I_{r+1}(N,q).
\end{equation}
Indeed each integer summand of the positive decreasing majorant
is at most its integral over $[n-1,n]$. For $r=0$ this also
bounds the tail in (CPP2). For the entire fixed-$k$ prime
contribution to $c_k$, (CPP1)--(CPP3) give the explicit bound
\begin{equation}\tag{CPP10}\begin{gathered}
 \left|\text{prime contribution to }c_k\text{ from }p^l>N\right|
 \le\frac14\sum_{r=1}^k\binom{k-1}{r-1}
                       \frac{b^{r-1}}{r!}I_{r+1}(N,q)
 \\\left(N\in\mathbb N,\ N\ge\exp((k+1)/q)\right).\end{gathered}
\end{equation}
This bounds the tail of the complete finite polynomial weight,
including all its signs. It is a bound for fixed $k,N,b$,
not a claimed uniform-in-$k$ cancellation estimate. For the
gamma tail after an integer $n=M\ge1$, the modulus of each factor in
(CPP6) is at most one, and
$\sum_{n>M}(q+2n)^{-2}\le1/[2(q+2M)]$.
Thus (CPP5) permits certified finite evaluations with complete
and explicitly bounded omissions.

\subsection{The original arithmetic prime points and the full support lift}

For the original arbitrary nontrivial bounded distributive
lattice $L$, the amplitude map
$\pi:G_L(\mathbb Z)\to\mathbb Z$, $(a,\lambda)\mapsto a$,
preserves addition, multiplication, and the unit by the defining
join/meet operations. Each ordinary prime $p$ in (CPP2)--(CPP10)
therefore has the exact prime preimage
\begin{equation}\tag{CPP11}
 \mathfrak p_p=\pi^{-1}(p\mathbb Z)
   =\{(0,\lambda):\lambda\in L\}
      \cup\{(a,1_L):a\in p\mathbb Z\setminus\{0\}\}.
\end{equation}
It is prime because $p\mid ab$ implies $p\mid a$ or $p\mid b$,
and it is proper because the unit is not in it. The supported
zero ideal is $(e)=\{(0,\lambda):\lambda\in L\}$: multiplication
of $e$ by $z_\lambda$ gives $z_\lambda$, and every such product
has zero amplitude. It is itself prime since $\mathbb Z$ has
no zero divisors, and $(e)\subsetneq\mathfrak p_p$. These
statements retain the supported zero stratum. They do not say
that (CPP11) exhausts the full split-support spectrum, or
assign an unproved numerical norm to the other lattice strata.
The integer $p$ and factor $p^{-lq}$ in (CPP2) are the original
Euler factors attached to these arithmetic prime preimages.

To specify the infinite trace without requiring an infinite
join in $L$, use the Banach vector space $\ell^1(I)$ on the
actual index set $I=\{(p,l):p\text{ prime},\ l\ge1\}$.
The linear maps $S_N(a)=\sum_{p^l\le N}a_{p,l}$ and
$S(a)=\sum_{p,l}a_{p,l}$ are defined with their actual
coefficients, and $S_Na\to Sa$ by absolute summability.
Equations (CPP2)--(CPP5) supply elements of this vector space
by the proved convergence. Their full maps are exactly
\begin{equation}\tag{CPP12}
 S_N^\uparrow(a,\lambda)=(S_Na,\lambda),\qquad
 S^\uparrow(a,\lambda)=(Sa,\lambda),\qquad
 \Phi_{\mathbb C}S^\uparrow=(S\times I)\Phi_{\ell^1(I)}.
\end{equation}
All labels are retained throughout the amplitude limit;
no infinite lattice join is introduced. The identities follow
by addition, distributivity, and direct substitution in
$\Phi_V(v,\lambda)=(v,1_K\otimes\lambda)$; its injectivity
and label retraction are proved in (CTP32)--(CTP34).
The same construction on $\ell^1(\mathbb N)$ applies to the
gamma sum. Finite linear combinations, including both
endpoint terms, the signs in (CPP3), and $A_k-E_k+P_k$, have
these same lifts. Any zero amplitude resulting from their
exact cancellation keeps its original label, and at top
support remains $e$, rather than the absence element $\tau$.

\subsection{A terminating search on every off-real input, without an equality trap}

The preceding formulas provide rigorous finite coefficient
enclosures for every fixed rational $b>1$. To include the
constant $c_0$, put $t=q/2>0$ and recall the convergent
digamma series, obtained from the same Euler limit for $\Gamma$:
\begin{equation}\tag{CPP13}\begin{gathered}
 \psi(t)=-\gamma+\sum_{n=0}^\infty
          \left(\frac1{n+1}-\frac1{n+t}\right),\\
 \left|\sum_{n>M}\left(\frac1{n+1}-\frac1{n+t}\right)\right|
       \le\frac{|t-1|}{M+\min(1,t)}
 \quad(M\in\mathbb N,\ M\ge1).\end{gathered}
\end{equation}
The summand equals $(t-1)/((n+1)(n+t))$; bound its absolute
value by $|t-1|/(n+\min(1,t))^2$ and integrate the decreasing
majorant from $M$ to infinity. Euler's constant itself has
the explicit enclosure, with $h_N=\sum_{j=1}^N1/j$,
\begin{equation}\tag{CPP14}
 h_N-\log N-\frac1N\le\gamma\le h_N-\log N
 \quad(N\in\mathbb N,\ N\ge1).
\end{equation}
Indeed $a_N=h_N-\log N$ decreases to $\gamma$ and
$a_n-a_{n+1}=\log(1+1/n)-1/(n+1)$ lies between zero and
$1/n-1/(n+1)$, by the integral of $1/x$ on $[n,n+1]$.
The telescoping tail is therefore at most $1/N$.
Together (CPP2), (CPP8)--(CPP10), (CPP13)--(CPP14), and the
gamma tail after (CPP10) give arbitrarily narrow rational
intervals for $c_0,c_1,\ldots,c_k$. Elementary constants can
also be enclosed: for positive rational $x$ use
$\log x=2\sum_{j\ge0}y^{2j+1}/(2j+1)$,
$y=(x-1)/(x+1)$, with its geometric remainder bound;
use the arctangent alternating series in Machin's identity
$\pi=16\arctan(1/5)-4\arctan(1/239)$ to enclose $\pi$,
whose tangent-addition formula and quadrant give the stated
identity, and interval monotonicity to enclose $\log\pi$.
Positive rational powers of the positive rational bases
can be enclosed by rational bisection of their defining
integer-power equations: for $q=A/B$ the relevant $B$th root
is the unique positive root, and rational bracket endpoints
are checked by exact integer powers. This includes both
$p^{-lq}$ and $N^{1-q}$. All truncation
errors tend to zero for fixed $k$ and $b$.

An explicit search now proceeds in finite stages $j=1,2,\ldots$.
At stage $j$ enclose $c_0$ and all $c_k$, $1\le k\le j$,
to width at most $2^{-j}$. Stop on the first pair of intervals
that proves $|c_k|>c_0$, recording the sign of $c_k$ and the
integer coefficient test $(1,-\operatorname{sgn}c_k)$ in
positions $0,k$. The inequality certificate and (CTP23)
prove negativity of the entire original trace. If an off-real
zero exists, (CTP22) supplies a finite $k$ with a strictly
positive gap $|c_k|-c_0$; for all sufficiently large stages
containing that $k$, the shrinking intervals prove the gap.
The search therefore terminates on every such zero configuration.
At no step does it wait forever for the sign of a single
possibly zero gap: all finite indices are revisited by the
stated stages. The real-root case has no such certificate.
This proves the computational content of the detector without
asserting that a negative coefficient has been found for the
actual original $H$.
